\section{Rebuttal}
Thank you very much for all the insightful reviews and suggestions. We have the questions listed and answered as follows.

For Reviewer 1:

\textbf{Q1: The proposed error model does make assumptions about the hardware fault models, which is acceptable since it already covers a large number of possible hardware failures, but it would be better if the paper states the coverage / limitation of hardware fault models more explicitly. For single error, the proposed model (that error perturbations follow a normal distribution) is only validated in the paper when there is one bit-flip in one layer output neuron. For multiple errors, the paper assumes each error is a random variable that follow a normal distribution, which means that every single error initially introduces one perturbation in one neuron value of a layer output. This assumption is valid for many memory errors. If we only assume one neuron is perturbed in the layer where failure occurs, it makes sense that single perturbation convolves with a random variable that has a normal distribution (i.e. next layer’s weights) generates perturbations that follow normal distribution. However, it might not be the case if the failure occurs internally in the logic portion of the accelerator, in which case one single error can generate multiple or even a large number of faulty neuron values in the layer output. The error only affect the neurons that use the faulty sequential logic, which is decided by the internal hardware logic instead of the convolution math.}

Ans1: This is a very good question and it is timely for the reliability study of neural network processing. In general, we utilize a random bit flip error model and inject random errors to weights and neurons in neural networks, which have been utilized in many prior fault analysis and fault injection works targeting at reliability issues about neural network processing as mentioned in this paper. Unfortunately, there is still a lack of systematic discussion about it. From our perspective, the bit flip error of neurons and weights is actually NOT directly correlated with any specific hardware faults like on-chip buffer errors, internal logic errors, and register errors. For instance, random bit flip errors in an on-chip buffer of a neural network accelerator can affect a neuron or a weight. Nevertheless, when the affected data is read and sent to the computing array like a systolic array, the data only affects a tiled computing task that utilizes this data rather than all the computing tasks that utilize this data. In addition, the same weight or neuron may also be read multiple times depending on the data flow utilized in the accelerator. As a result, transient faults in the buffer may probably happen on only one of the read operation rather than all the read operations. Essentially, the influence of realistic faults can depend on the underlying microarchitecture, data flow, and compiler of the neural network accelerator, which is more complex than the bit flip fault model. As for the internal logic errors, SEU of input registers and accumulator registers in neural network accelerators may partly show as a bit flip error in neurons or weights. Basically, we want to argue that the bit flip error model is not a realistic hardware fault model. Instead, we would like to take it as an abstract application-specific fault model for neural network processing despite the underlying computing architectures which can either be neural network accelerators, GPUs, and CPUs. It is somehow similar to program fault analysis that relies on instruction-wise fault injection[2][3] on CPUs and GPUs. In general, the advantage is that it is convenient to model the influence of soft errors on neural network processing without sticking to any specific architectures and compilers. The disadvantage is the limited accuracy to specific architectures. The accuracy of the model is partly verified in [1] with a voltage down scaling experiment, but there is a lack of systematic investigation about it. This remains a valuable problem and we may further investigate it in future. 

[1] Reagen, Brandon, Udit Gupta, Lillian Pentecost, Paul Whatmough, Sae Kyu Lee, Niamh Mulholland, David Brooks, and Gu-Yeon Wei. "Ares: A framework for quantifying the resilience of deep neural networks." In 2018 55th ACM/ESDA/IEEE Design Automation Conference (DAC), pp. 1-6. IEEE, 2018.

[2] Sridharan, Vilas, and David R. Kaeli. "Eliminating microarchitectural dependency from architectural vulnerability." In 2009 IEEE 15th International Symposium on High Performance Computer Architecture, pp. 117-128. IEEE, 2009.

[3] T. Tsai, S. K. S. Hari, M. Sullivan, O. Villa and S. W. Keckler, "NVBitFI: Dynamic Fault Injection for GPUs," 2021 51st Annual IEEE/IFIP International Conference on Dependable Systems and Networks (DSN), Taipei, Taiwan, 2021, pp. 284-291, doi: 10.1109/DSN48987.2021.00041.

\textbf{Q2: Moreover, the values of these faulty neuron values are correlated and they can not be seen as independent variables. It seems that current statistic models do not cover such cases.}

Ans2: As for the correlation between different errors of neurons, we agree that the neurons are not independent as they are correlated from both the perspective of underlying circuits and neural network architecture that connects the neurons. The number of correlated neurons and the connection between these neurons are extremely large and complex. In addition, the numerical errors rapidly spread across almost all neurons in the following layers according to the typical neural network architectures. Hence, it is rather challenge to model the correlation from the perspective of a single neuron directly. In this paper, the identical independent neuron assumption, however, is utilized to simplify the modeling and approximate the reliability of the overall neural network rather than a specific neuron. As a result, although the analysis is demonstrated to be good enough for the overall neural network model accuracy, it CANNOT be utilized to characterize the difference between neurons. A possible analogy may be Brownian motion that can utilize random particle movement model to characterize the behavior of gas or liquid from a macro perspective. At the same time, the analysis does NOT deny the correlated force acting between particles. Basically, we may regard indicators such as RMSE as the energy or temperature of particles, and error propagation as the collision and energy transfer between a large number of particles, then we can estimate the whole system from a macro perspective without knowing all the microscopic states. 

\textbf{Q3: The proposed error model do not consider layers other than MAC(Multiply Accumulate) layers. One example is that BatchNorm layers seem to play an effect in Fig 5, where Resnet18 demonstrates a more stable RRMSE than VGGNet. I suppose Resnet18 has BatchNorm while RRMSE does not. However, it doesn’t seem that the introduction of BatchNorm layer will complicate the entire modeling process that much, since normalization will only make the network more ideally distributed as normal distributions. It seems that normalization layers such as BatchNorm will help the existing error model. Another example is activation layers. Lemma 2 and lemma 4 might need to be adjusted when ReLU is introduced. But ReLU won’t affect the distribution too much as evidenced by the experiment results. It would be helpful if the author can present the assumptions more explicitly, and discuss how the introduction of other layers such as normalization and activation can affect the modeling, and more importantly, why the experiment results still match the real cases even though these layers were not considered throughout the modeling.}

Ans3: The batchnorm layer is essentially a linear operator that scales the error and activation equally, so it does not change RRMSE theoretically. Your observation about the influence of batchnorm is correct. Batchnorm layer can make the distribution of weights and activation values closer to zero-mean normal distribution hypothesis(\autoref{eq:normal_hyp}), which makes the statistical model conform better.

As for the non-linear activation functions, they will affect the outputs directly. For example, ReLU will filter out computing errors when activation is negative. Suppose the propagated errors distribute evenly on positive and negative activation values, it can be expected that ReLU will filter out half activations and errors simultaneously. Then, Lemma 2 and lemma 3 can approximately change to $\mathrm{var}(x_{l+1})\approx \frac{m}{2}\mathrm{var}(x_l)\mathrm{var}(w_l)$ and $\mathrm{var}(\Delta_{l+1})\approx \frac{m}{2}\mathrm{var}(\Delta_l)\mathrm{var}(w_l)$ if we merge the affect of ReLU in to it. Thereby, RRMSE remains unchanged. In summary, for typical neural networks built with convolution-batchnorm-ReLU layers, neural network depth will not affect $RRMSE$ i.e. neural network resilience in general. However, in practice, the above assumption can be different. For instance, the percentage of positive neurons and negative neurons can be unequal, which may affect the RRMSE estimation accordingly. 

We add the above description to Section IV.A and cover the typical operators in neural networks in the fault analysis. However, we did not modify Lemma in III.B because the influence of activation function is not considered here. In addition, there are many different kinds of activation functions other than ReLU, and precise modeling and experiments are too complex, so we only briefly discuss the influence of commonly used ReLU function in Section IV.A.

For Reviewer 2:

\textbf{Q1: In figure 1, the activations do not fit well with a normal distribution because a skew is observable in the graphs. They seem to fit better with a lognormal distribution. Please clarify.} 

Ans1: In order to quantitatively measure the accuracy of normal distribution model, we add three classical statistical parameters in Fig.1 including mean, skewness, and excess kurtosis. It can be observed that mean is generally close to zero and skewness is generally less than 0.5, which confirms the fitness of normal distribution model to the weights and activations. On top of the modeling accuracy, the nice statistical features of normal distribution model such as aggregation also make it easier to understand and analyze in depth. Hence, it is a competitive candidate for modeling and adopted in this work. 

We agree that it is possible to replace normal distribution with more accurate models. We tried lognormal and t-LOC respectively. Since lognormal requires data to be greater than 0, additional location parameter is required and it includes three parameters in total. However, we find that the log-likelihood value of 3-parameter lognormal is usually worse than the 2-parameter normal distribution, and the lognormal parameters obtained with maximum likelihood estimation are unstable. For t-LOC, it fits the neurons better than normal distribution on some data. However, the Akaike information criterion(AIC)[1] that indicates the model generality is worse due to more parameters. In Fig.1, 5/9 of the data fits better with t-LOC compared to normal distribution, but only 2/9 of the data shows better AIC. Basically, the advantage of t-LOC compared to normal distribution model is trivial while it poses more challenges on in-depth analysis of the overall neural network reliability.

[1] Akaike, H.: A new look at the statistical model identification. IEEE Trans. Autom. Control 19(6), 716–723 (1974)

\textbf{
Q2: The authors assumed that "the distribution of the weights and activations are independently and identically to simplify the modeling in the rest of this work." This assumption is clearly wrong because of the impacts of reconvergent paths. An error in one weight can be propagated through different paths and then recovergent in the next layers. }

Ans2: Similar question is also raised by Reviewer 1, please refer to Ans2 for Reviewer1.

\textbf{Q3: It is required to re-draw the results of the experiments on page 8 with a less range for RRMSE. It can better show the gap between the FI and empirical results when the range of X-axis decreases. Also, the accuracy is not directly compared with RRMSE. Is there any other way to compare these two metrics directly so that we can be sure about their closeness?} 

Ans3: The original figure is a bit misleading, we updated the Y-axis of Fig.3 in Page 8 such that it reveals the correlation between RRMSE and model accuracy more clearly.

As for the difference between FI and empirical results, we add "Accuracy Difference" on right Y-Axis to highlight the quality of the empirical model. 

We are not quite sure with the question "Is there any other way to compare these two metrics directly?" Fig.3 directly shows the correlation between RRMSE and accuracy based on both models and FI experiments. 

\textbf{Q4: How does Figure 4 show the fitness of the real RRMSE and the predicted one? The slopes of the lines are not the same! I recommend to compare the more common ways to compare these two value that is to show either the absolute difference or relative difference between them!}

Ans4: We add diagonal which is the ideal correlation between real RRMSE and the predicted RRMSE in the figure such that the difference is presented more clearly. At the same time, we also add "Relative Difference" on the right Y-Axis and quantize the RRMSE difference. Note that Relative Difference = (Predicted RRMSE - Real RRMSE)/Real RRMSE. From Fig.4, we can observe that the RRMSE difference is more significant when RRMSE gets larger in general, which may be caused by the error correlation ignored in the model at higher error rate.

Reviewer 3:

\textbf{Q1: A traditional and accurate method for assessing DNNs' reliability has been resorting to fault injection, which suffers from prohibitive time complexity. In comparison, analytical and hybrid analytical/fault injection-based methods have been proposed. This paper lacks comparison with the state-of-the-art works for reliability assessment of DNNs based on the analytical models. Analytical models are accurate and can be applied to various ranges of NN architectures. Yet, they do not need to come with fault injection and can be used independently for reliability analysis. Therefore, a motivation on using a statistical model instead of other reliability analysis methods is a must. The imbalance between fault injection methods and other approaches (analytical, hybrid) is very noticeable in the related work. The authors should provide a better justification. If more studies are available in these domains (analytical, hybrid), the authors should enhance the related work section with the papers from these sub-domains.}

Ans1: There have been efforts made to understand artificial neural networks from statistical perspectives [1] [2] and characterize soft error rate with statistical models [3] [4], which demonstrated the potential of using statistical models for neural networks and soft errors. We also notice that BinFi [5] proposed to take advantage of an analytical model to simplify the the fault injection through a binary search algorithm. However, the model is more like a qualitative conclusion. Basically, the authors argue that errors inducing larger data variation generally pose more significant influence on model accuracy compared to the errors leading to smaller data variation. There is a lack of in-depth analysis with models. 

We add the above description to related work as well. However, there is still a lack of analytical work on investigating the influence of soft errors on convolutional neural networks in general at the moment. 

[1] Cheng, Bing, and D. Michael Titterington. "Neural networks: A review from a statistical perspective." Statistical science (1994): 2-30.

[2] White, Halbert. "Learning in artificial neural networks: A statistical perspective." Neural computation 1, no. 4 (1989): 425-464.

[3] Chang, Austin C-C., Ryan H-M. Huang, and Charles H-P. Wen. "CASSER: A closed-form analysis framework for statistical soft error rate." IEEE Transactions on Very Large Scale Integration (VLSI) Systems 21, no. 10 (2012): 1837-1848.

[4] Peng, Huan-Kai, Charles H-P. Wen, and Jayanta Bhadra. "On soft error rate analysis of scaled CMOS designs: a statistical perspective." In Proceedings of the 2009 International Conference on Computer-Aided Design, pp. 157-163. 2009.

[5] Chen, Zitao, Guanpeng Li, Karthik Pattabiraman, and Nathan DeBardeleben. "Binfi: An efficient fault injector for safety-critical machine learning systems." In Proceedings of the International Conference for High Performance Computing, Networking, Storage and Analysis, pp. 1-23. 2019.

\textbf{Q2: In the third paragraph of the introduction section, it is said that some simulation results may lead to contradictory conclusions. Referring to the paper entitled: “Statistical fault injection: Quantified error and confidence”, to find the required number of repetitions for the fault simulation experiments, an equation is provided that can reach 95\% confidence level and 1\% error margin in fault injection results. Therefore, this statement is not true if a sufficient number of faults are selected for the fault injection process. I recommend authors elaborate further on this sentence.}

Ans2: There is some misunderstanding here. It is not the number of fault injections that leads to the error confidence and contradictory conclusions. In fact, both experiments in Ares[1] and SNR[2] have sufficient number of sampling to obtain bit error rate-accuracy curve with sufficient confidence. However, they have distinct model quantization setups that result in different conclusions. The bit error rate-accuracy curve obtained is only valid for a specific quantization setup. 

[1] Reagen, Brandon, Udit Gupta, Lillian Pentecost, Paul Whatmough, Sae Kyu Lee, Niamh Mulholland, David Brooks, and Gu-Yeon Wei. "Ares: A framework for quantifying the resilience of deep neural networks." In 2018 55th ACM/ESDA/IEEE Design Automation Conference (DAC), pp. 1-6. IEEE, 2018.

[2] Ozen, Elbruz, and Alex Orailoglu. "SNR: Squeezing Numerical Range Defuses Bit Error Vulnerability Surface in Deep Neural Networks." ACM Transactions on Embedded Computing Systems (TECS) 20, no. 5s (2021): 1-25.

\textbf{Q3: In subsection (A) of section IV, the authors do not elaborate on the reason they set the bit error rate to be 1E-6. The fault model and fault injection can be also discussed.}

Ans3: We improve the description about the bit error rate setup in Section IV. In general, the bit error setups are not realistic data from specific application scenarios. Instead, they cover a broad range of error scenarios that lead to different model accuracy drop. 1E-6 is selected in the experiment because it stands for a typical scenario that results in notable accuracy drop. %In addition, we find that the results in Section IV.A are not sensitive to the bit error rate setting. Due to the error influence aggregation mentioned in section III.E, the relative RRMSE difference of the two layers in Fig. 5 will not change significantly under different bit error rate.

The description about fault injection is also improved in Section IV, the random bit flip error rate refers to faulty neuron bits over all the neurons in a neural network model which is consistent with prior fault injection frameworks including Ares [1] and PyTorchFI [2].

[1] Reagen, Brandon, Udit Gupta, Lillian Pentecost, Paul Whatmough, Sae Kyu Lee, Niamh Mulholland, David Brooks, and Gu-Yeon Wei. "Ares: A framework for quantifying the resilience of deep neural networks." In 2018 55th ACM/ESDA/IEEE Design Automation Conference (DAC), pp. 1-6. IEEE, 2018.

[2] Mahmoud, Abdulrahman, Neeraj Aggarwal, Alex Nobbe, Jose Rodrigo Sanchez Vicarte, Sarita V. Adve, Christopher W. Fletcher, Iuri Frosio, and Siva Kumar Sastry Hari. "Pytorchfi: A runtime perturbation tool for dnns." In 2020 50th Annual IEEE/IFIP International Conference on Dependable Systems and Networks Workshops (DSN-W), pp. 25-31. IEEE, 2020.

\textbf{Q4: In subsection (A) of section (V), what type of quantization is applied to VGGNet-11 and ResNet-18? Are these fully quantized networks including internal layers’ data and activations of the network?}

Ans4: The quantization setup is the same with that described in section IV. Specifically, it is a simplified post-training quantization for each layer because quantization scale parameter is set after training. The range of the quantizer is symmetric and determined by the maximum range of input values.

\textbf{Q5: An informative table is expected to compare the accuracy results of this method and other vulnerability assessment methods and the FI process.}

Ans5: To the best of our knowledge, neural network vulnerability is mainly obtained with fault injection and realistic measurement under radiation test. It is difficult to conduct radiation test for us. As for simulation based fault injection, we follow the same fault injection strategy proposed in Ares [1] and PyTorchFI [2] and re-implemented it in PyTorch, so fault injection is already compared with modeling based method in this work. The main reason that we did not use PyTorchFI directly is the limited quantization support that cannot fit the experiment requirements in this work.  

[1] Reagen, Brandon, Udit Gupta, Lillian Pentecost, Paul Whatmough, Sae Kyu Lee, Niamh Mulholland, David Brooks, and Gu-Yeon Wei. "Ares: A framework for quantifying the resilience of deep neural networks." In 2018 55th ACM/ESDA/IEEE Design Automation Conference (DAC), pp. 1-6. IEEE, 2018.

[2] Mahmoud, Abdulrahman, Neeraj Aggarwal, Alex Nobbe, Jose Rodrigo Sanchez Vicarte, Sarita V. Adve, Christopher W. Fletcher, Iuri Frosio, and Siva Kumar Sastry Hari. "Pytorchfi: A runtime perturbation tool for dnns." In 2020 50th Annual IEEE/IFIP International Conference on Dependable Systems and Networks Workshops (DSN-W), pp. 25-31. IEEE, 2020.

\textbf{Q6: Comparing the processing time and generalization of this approach with other state-of-the-art methods needs to be addressed}

Ans6: There is still no prior statistical modeling based neural network reliability analysis work. Currently, we take fault injection based approach as the reference and verify the proposed modeling accuracy.

\textbf{Q7: Is it further extendable to extract neurons’ vulnerability or neuron-wise vulnerability ranges expressing whether a fault would misclassify the network or not? It will enable a comprehensive reliability study where so-called vulnerability factors of layers, neurons, and bits in a DNN can be obtained.}

Ans7: This is a very good question. Basically, the proposed statistical model is mainly valid for reliability analysis of the entire neural network that includes a large number of neurons, but it cannot be applied to reliability analysis for specific neurons because the model assumption does not distinguish the neurons at the very beginning. More detailed discussion can also be found in the answer to Q2 from Reviewer 1. 